\begin{recipe}{Horiatiki}
\label{recipe:horiatiki}
\kicker[Bijgerecht,Vegetarisch,Grieks]{Bijgerecht · vegetarisch · Grieks}
\index[register]{Horiatiki}
\index[register]{Feta}
\index[register]{Tomaten}
\index[register]{Komkommer}
\index[register]{Olijven}

\meta{10 minuten}

\lettrine{H}{oriatiki} is de Griekse boerensalade: tomaat, komkommer, rode ui en
olijven met een hele plak feta erbovenop, alleen olijfolie en oregano erdoor.
Geen sla en geen dressing, dat is geen vergissing maar traditie.

\blockrule{Ingrediënten}
\begin{ingredients}
  \ing{200 g}{feta, in één plak}\index[register]{Feta}
  \ing{4}{trostomaten, in partjes}\index[register]{Tomaten}
  \ing{1}{komkommer, in halve manen}\index[register]{Komkommer}
  \ing{1}{rode ui, in dunne ringen}
  \ing{100 g}{zwarte olijven zonder pit}\index[register]{Olijven}
  \ing{3 el}{olijfolie}
  \ing{1 tl}{gedroogde oregano}
  \ingb{zout, naar smaak}
\end{ingredients}

\blockrule{Bereiding}
\tip{Het klassieke bijgerecht bij moussaka: fris-zuur tegen de bechamel.}
\begin{steps}
  \item Was de tomaten en de komkommer. Snijd de tomaten in partjes en de
        komkommer in halve manen. Snijd de rode ui in dunne ringen.
  \item Verdeel tomaat, komkommer, ui en olijven over een schaal of bord.
  \item Leg de plak feta boven op de salade in zijn geheel, niet verkruimeld.
        Besprenkel met de olijfolie, bestrooi met oregano en een beetje zout.
        Serveer direct.
\end{steps}

\heroimagefade[fadewidth=0.8cm]{images/horiatiki}
\end{recipe}
